% SPDX-License-Identifier: Apache-2.0
\section{What the Runtime Is}
\label{sec:r-overview}

The embedding article~\cite{embedding} argues that a hardware
description language embedded in Rust needs no parser and no keywords,
and that what remains is a library. The crate \code{txhdl} is that
library, and this document is the crate. Every file below is included
from the tree when the document is built, so a listing here cannot
differ from the code, and the rule the project keeps is that a concept
enters the runtime, a compiled example and these documents in the same
change, or the change is not finished.

The crate is small on purpose. Four modules and one procedural macro
crate hold everything:

\begin{itemize}
\item \code{txhdl::types}: two-valued \code{Bit}, nine-valued
  \code{Logic}, the numeric vectors \code{U<N>} and \code{I<N>},
  \code{logic::Vec<N>}, and the two marker traits \code{Transaction}
  and \code{Tag}.
\item \code{txhdl::comp}: clocks, wires and channels and their ends,
  interfaces, crossings, registers and memories, units, configurations,
  and the executor.
\item \code{txhdl::pipeline}: operators that take a cycle, all
  \code{async}.
\item \code{txhdl::funcs}: operators that cannot, all plain \code{fn}.
\item \code{txhdl\_macros}: the two derives, \code{interface!},
  \code{when!} and \code{case!}, without \code{syn} or \code{quote}.
\end{itemize}

The whole crate is a simulation-shaped prototype: values exist while it
runs, wires are shared cells, and a testbench may look at anything.
The lowering to hardware is the experiment the article states, and this
crate is what that experiment would be written against.

\subsection{The model of time}
\label{sec:r-time}

The executor is the one part of the runtime a reader cannot recover
from the signatures, so it is stated here before the code.

\textbf{Time runs in ticks.} One poll of the top unit is one tick.
Every clock is a type with a period and a phase in ticks, and has a
rising edge at every tick \code{t} with \code{t >= PHASE} and
\code{(t - PHASE) \% PERIOD == 0}, and a falling edge half a period
later. The default clock has period 2 and phase 0, so its cycle is two
ticks and its falling edge is a tick of its own; a single-clock design
mentions neither number, and its testbench steps by
\code{Running::cycle}.

\textbf{A process waits, then reads.} The waits are events:
\code{C::rising()} and \code{C::falling()}, the next edge of a clock;
\code{rx.wait()}, the next rising edge of the channel's clock at which
a transaction is offered, taken; and \code{until(C::rising, || cond)},
or the same with \code{falling}, the next such edge at which a
condition on state holds. Reads are plain: \code{Reg::get} returns
what the last edge latched, a wire's \code{get} returns what it
carries, a memory's \code{read} is its asynchronous port. Drives are
deferred: \code{Reg::set} and \code{Mem::write} take effect at the end
of the step, so a read after a drive in the same iteration still sees
what the edge latched.

\textbf{Every \code{.await} is a cycle boundary.} A process that has
crossed an edge at this step cannot cross it again, so a second wait
written after a first waits for the next edge, as an operator does.
Waits that share one edge are written inside \code{parallel!}, which
polls its futures together and returns their outputs as a tuple; that
is the one case the executor treats specially. Once one wait of the
group has crossed the edge at this step, the others in the group are
free, each once, so the group is one wait. State read before a wait,
at the top of a loop, is read in the previous step and before that
step's drives commit; state is read after the wait.

\textbf{A process is its waker.} The executor tells processes apart by
the data pointer of the waker they are polled with. \code{join2} and
\code{join\_all} give each child a fresh one, so the rule above holds
per process rather than per unit. \code{parallel!} keeps its futures in
the one process, which is the difference between the two: \code{join2}
forks, \code{parallel!} does not.

\textbf{An operator's cycle is in the process's clock.} An operator in
\code{txhdl::pipeline} waits one edge of whatever clock the process
last crossed an edge of, which is the only sense a latency has inside
a process. A process that has crossed no edge yet is in the default
clock.

\textbf{What the executor does not check.} A process that reads a
register of another clock without a \code{Crossing} is a hazard, and
the prototype runs it. The register's type says its clock; the
process has no type that says one. The asynchronous FIFO example
commits this hazard on purpose, on its memory, under the pointer
discipline that makes it safe.

\subsection{What is prototype only}

Two things in the code exist for the simulation and would not be
lowered. \code{Running} and \code{simulate} drive a design from a
\code{main}, and \code{now} reports the time step for traces. And the
shared cells behind every wire, channel and register, with the deferred
drives that commit at the end of a step, are how a simulation shares
state; the lowering would replace them with nets and flops.

\subsection{Tracing}

A design can be watched as a waveform. \code{comp::trace} holds a
registry of probes, each a dotted path, a width and a closure that
reads the signal; \code{\#[derive(Trace)]} on a unit registers every
field under its field name, so the hierarchy is the nesting of units
and nothing is named at construction. Registers, wire ends, channel
ends and crossings register themselves, plain values register nothing,
and \code{\#[derive(Value)]} gives a struct or a fieldless enum of the
design's own its width and bits. \code{Vcd} is the sink: the executor
calls it at the end of every tick, after the drives have committed,
and it writes what changed; a clock is written as the level it is,
high from its rising edge to its falling one. FST would be the same probes with another
writer, and there is one now: \code{Fst}, over the \code{fst-writer}
crate, the project's first external dependency, pinned by lockfile
through \code{crate\_universe}. \code{Wave::from\_env} chooses the sink
by the environment, \code{TXHDL\_FST} or \code{TXHDL\_VCD}, and
\code{stop} finishes the file. A compound value is one bit vector and
one signal per field beside it, from \code{Value::parts}; the writer
has no string variables, so a value keeps its bits.
The examples document draws a VCD through \code{vcdcvt} and
\code{sqlite2drawtiming}, prebuilt tools pinned by checksum, and
\code{//tools/dt2tikz}, which turns drawtiming's text into TikZ.

\subsection{Netlists}

The probe registry records, for every signal, its kind and the shared
cell it is on, so the two ends of one wire match. \code{netlist::verilog}
walks a design through it and writes a Verilog skeleton: a module per
unit, a \code{reg} per register, ports where a unit holds one end of
a wire whose other end is elsewhere, and instances with their ports
connected. It sees fields only, so a unit whose ports are parameters
of \code{run} yields its registers and nothing else, and it writes no
behaviour. That is the structural half of the lowering. The
behavioural half has begun: \code{\#[pipeline]} lowers an \code{async
fn} of awaited operators, and \code{\#[lower]} lowers a unit's
\code{run} loop, with \code{netlist::unit\_verilog} assembling the
module from the unit's fields, the ports of \code{run} and the
statements the macro read, \code{case!} and channels included. The
macro builds an expression tree, \code{netlist::Expr}, and a
\code{Lowered} unit around it, and two emitters render it, Verilog
and VHDL-2008. The VHDL is what closes the loop: the build simulates
a lowered unit under \code{nvc}, with a testbench
\code{//tools/fst2tb} generates from the trace the Rust run wrote, and
the test passes only if registers and outputs agree at every tick.
What neither emitter reads yet is a unit's second process, a
falling-edge wait, and a memory; and a unit with channels cannot be
checked until the runtime's channel is a cycle-accurate handshake
rather than a mailbox.

\subsection{Building it}

\begin{lstlisting}[style=ascii,caption={The runtime and everything written against it, built hermetically.},label={lst:r-build}]
bazel build //lib/...                 # the crate and the macros
bazel build //lib/examples/...        # every example
bazel run //lib/examples:ex_fifo      # one that prints a trace
bazel run @rules_rust//:rustfmt \
  --@rules_rust//:rustfmt.toml=//:rustfmt.toml -- //lib/...
\end{lstlisting}

The last line formats the tree at 80 columns, which is the width these
documents include a file at. A line longer than that overflows its
frame, visibly, rather than wrapping.
